\lyseis{6.2}{Λόγος και αναλογία}
\sol{1} α) $\frac{3}{4}$ \quad β) $\frac{30}{120}=\frac{1}{4}$ \quad γ) $\frac{250}{50}=5$ \quad δ) $\frac{0{,}6}{1{,}5}=\frac{2}{5}$ \quad ε) $\frac{3}{4}:\frac{1}{2}=\frac{3}{2}$.
\sol{2} α) Ναι. β) Όχι ($4\cdot9=36\neq35=7\cdot5$). γ) Ναι ($12\cdot30=360=18\cdot20$). δ) Ναι ($0{,}5\cdot12=6=2\cdot3$). ε) Όχι ($7\cdot25=175\neq168=8\cdot21$).
\sol{3} α) $15x=60$, $x=4$ \quad β) $12x=72$, $x=6$ \quad γ) $5x=60$, $x=12$ \quad δ) $21x=84$, $x=4$ \quad ε) $5x=15$, $x=3$.
\sol{4} $4\cdot15=60=6\cdot10$. Οι δύο λόγοι είναι ίσοι ($\frac{2}{3}$ ο καθένας).
\sol{5} $\frac{18}{27}=\frac{2}{3}$.
\sol{6} α) $4\cdot50\,000$ cm $=200\,000$ cm $=2$ km. β) $15$ km $=1\,500\,000$ cm, διά $200\,000$ δίνει $7{,}5$ cm. γ) $8\cdot1\,000$ cm $=80$ μ.
\sol{7} α) Σωστό. β) Λάθος: $\frac{200}{50}=4$. γ) Σωστό. δ) Σωστό. ε) Λάθος: έχουν ίσους λόγους περιμέτρων, όχι ίσες περιμέτρους.
\sol{8} $3{,}6\cdot5=18$ km και $5{,}2\cdot5=26$ km.
\sol{9} $42$ km $=4\,200\,000$ cm, διά $200\,000$ είναι $21$ cm.
\sol{10} Λόγος μεγέθυνσης $\frac{40}{10}=4$, άρα νέο ύψος $4\cdot15=60$ cm. Περίμετροι: $2(10+15)=50$ και $2(40+60)=200$, λόγος $\frac{200}{50}=4$.
\sol{11} Το $3\times1{,}2$ μ.: $300:25=12=120:10$.
\sol{12} Λόγος $1:4$, άρα σιρόπι $750:4=187{,}5$ ml.
\sol{13} $36\cdot25=900$ cm $=9$ μ.
\sol{14} $\frac{15}{x}=\frac{3}{2}$, $3x=30$, $x=10$ λίτρα.
\sol{15} Το πλάτος γίνεται $2$ φορές μεγαλύτερο ($\frac{2}{1}$) αλλά το ύψος μένει ίδιο ($\frac{1}{1}$): οι λόγοι δεν είναι ίσοι.
\sol{16} $\frac{4}{10}=0{,}4$ και $\frac{6}{15}=0{,}4$: όμοια. Για τα $4\times6$ και $8\times10$: $\frac{4}{8}=0{,}5$ και $\frac{6}{10}=0{,}6$: όχι όμοια.
\sol{17} Πολλαπλασιάζω και τα δύο μέλη με $\beta\delta$: $\frac{\alpha}{\beta}\cdot\beta\delta=\alpha\delta$ και $\frac{\gamma}{\delta}\cdot\beta\delta=\beta\gamma$, άρα $\alpha\delta=\beta\gamma$.
\sol{18} $\frac{12}{16}=\frac{3}{4}$ και $\frac{6}{8}=\frac{3}{4}$.
\sol{19} $7{,}5\cdot10$ km $=75$ km. Ο χάρτης $1:200\,000$ δείχνει περισσότερες λεπτομέρειες, γιατί η ίδια απόσταση σχεδιάζεται μεγαλύτερη.
\sol{20} Ελεύθερη απάντηση.
\sol{21} Ελεύθερη απάντηση. Οι περισσότερες οθόνες έχουν λόγο μήκους προς πλάτος $\frac{16}{9}$ ή παραπλήσιο.
